\section{EAST 破裂预警的炮内和炮间更新}
\label{sec:east360-online-adaptation}

我们在 EAST 的历史放电记录上检验一个用 J-TEXT 数据训练的破裂预警模型。一次放电简称一炮。放电过程中，预警模型根据当时已经收到的诊断信号给出破裂风险分数。另一个模型负责预测接下来一小段诊断波形。实测波形到齐后，可以比较预测和实测，再修正本炮之后的风险分数，这就是炮内更新。炮间更新则要等一批放电结束，用其中少量已知结局调整以后各炮的评分。两种更新都保留 J-TEXT 预警模型的权重。

\subsection{怎样评价一次报警}

主实验按同一顺序回放 EAST 的 360 炮，其中 180 炮发生破裂，180 炮未发生破裂。风险分数第一次达到报警阈值时，记为该炮的首次报警。对于破裂炮，只有首报提前量大于 30 ms 才算及时，恰好 30 ms 不计入。对于未破裂炮，只要整炮中出现过一次报警，就记为一炮误报。因此，及时检出率和误报率的分母各为 180。主要比较看的是：在误报率不超过 20\% 时，能否让更多破裂炮及时报警。

前 20 炮用于准备 EAST 的评分方法，其中 8 炮破裂、12 炮未破裂。这 20 炮在各项对照中使用相同的预测和 J-TEXT 阈值 0.7258947，照常计入总成绩。保存并锁定预测后，才读取结局，拟合 EAST 的基础评分方法。确定报警阈值时，轮流留出一炮未破裂记录，用其他 19 炮拟合，再记录留出炮的最高风险分数，共做 12 次。主结果只比较一个 EAST 报警阈值 0.5062586，由这 12 次的结果确定，用于后 340 炮。后 340 炮的结局不用于重定阈值或重拟合基础评分。

\subsection{从诊断信号到风险分数：原有的四步修正}

先看静态模型怎么打分。每个时刻，模型读入截至当前的 128 个采样点，包括 23 路诊断信号以及哪些点缺测的信息。每 8 点分成一块，共 16 块。模型把每块转成一个向量，再用两层注意力汇总成一组 64 维特征。评分还用到 23 个统计量，描述各路信号近期的变化方向是否一致，以及一个小网络从当前窗口算出的 24 维时序特征。评分层给这三组数中的每一个乘以固定权重，把结果加起来，再换算成 0 至 1 的风险分数。静态模型本身已经能看一段历史，只是不拿此前的波形预测误差来修正评分。

\textbf{第 1 步：记下本炮的预测偏差。} 每当一段预测有了实测，就按通道算两件事：这一段整体预测偏高还是偏低，误差在这一段是往上走还是往下走。将这两个量和此前记下的值各取一半，得到新的记录。模型根据当时的等离子体电流、磁场和电流变化，用固定的映射把这份记录转成对 24 维时序特征的补充。编码器读到的信号不变。这样，即使两个时刻眼前的信号相近，评分也可以因本炮先前的预测表现不同而有区别。

\textbf{第 2 步：告诉评分层，最近一段偏离了多少。} 电流或磁场变了，平常的预测误差也可能不同。模型用本炮较早且工况相近的波形误差，结合 J-TEXT 未破裂炮的统计量，估计在该工况下通常会出现多大的误差。最近一段相对这个参考偏高还是偏低，决定修正方向；相邻两段是否都朝同一方向偏、偏得有多明显，则限制修正大小。这项信息也加到同一组 24 维特征上，第 1 步的记录仍保留。前一步记的是长期积累的偏差，这一步补的是最新误差比相似工况下的常见误差多了多少。

\textbf{第 3 步：调整编码器看到的信号走势。} 模型比较预测与实测中不同诊断信号的先后变化。例如，位置变化之后的辐射变化是否与预测相符。它用梯度估计哪些通道之间的调整能让预测更接近这种变化关系，再把这个方向换成修正系数。对输入窗口中的一个目标点，取另一组信号在此前第 4 点与第 12 点的差，据此对目标点作小幅调整。这个步骤在时间块编码前执行，保持各通道在窗口中的平均水平，不移动采样点，也不填补缺测。前两步的 24 维特征保留，变的是编码器从调整后的输入提取出的 64 维特征。

\textbf{第 4 步：再处理输入中的整体偏差。} 模型用本炮此前已经完成的预测估计各通道的平均误差，误差按实测减预测计算。在第 3 步调整后的输入窗口里，再从有实测的点扣除这个偏差估计，并限制改动大小。第 3 步改变窗口内的起伏，第 4 步允许调整信号的整体水平。保留这四步的版本，就是用来检验下一步的主要对照。

\subsection{新增炮内修正：等实测到齐，再调整时间块表示}

诊断波形预测每隔 32 个采样点发出一次。例如，在第 \(q\) 点发出的预测，目标是第 \(q+1\) 至 \(q+25\) 点的诊断信号。回放按 EAST 记录中的原始采样时刻推进，名义采样率为 1 kHz。整段实测波形最早在第 \(q+26\) 点可用于计算修正。主实验再等 5 个采样点，从第 \(q+31\) 点开始过渡，当点仍用旧修正，随后用 8 个采样间隔逐步换成新修正。等待期间预警模型仍按当时已有的信息给分，已给出的分数不会被改写。

\textbf{第 5 步：用新误差调整模型内部的时间块表示。} 将预测和实测各分成 5 小段，比较每段是预测偏高还是偏低。再扣除发出预测时就能从本炮历史估计的误差，留下这段实测带来的新误差。模型用梯度做两次换算：先判断原来的哪些内部特征会影响波形预测，再判断在目前的输入窗口里，改动哪些时间块能让这些特征朝误差提示的方向变化。磁扰动、位置、辐射和环电压四组信号分别参与计算。

这一步加在时间块已编码成向量之后、第一层注意力之前。举例说，若一对相邻时间块的向量是 \(u\) 和 \(v\)，一组信号给出的修正为 \(d\)，就改成 \(u-d\) 和 \(v+d\)。两块的总和不变，差异却变了，注意力层看到的前后变化也会不同。随后用同一套编码器和评分系数重算当前分数。改的是当前窗口的表示，原始信号和模型权重都保持不变。这些临时量到炮末清零，且不读取本炮是否破裂或破裂时刻。

主要对照是只做到第 4 步。另一个对照也做到第 5 步，但把刚到齐的误差换成发出本次预测时已经算出的上一笔合格误差。两种方法用同一套计算求修正方向，在相同的时刻和位置施加修正，并使每一时刻实际改动的大小相同。这样比较的是哪一笔实测误差提供了更有用的信息。

\subsection{炮间：先保存整批预测，再用已知结局调整后续评分}

后 340 炮每 20 炮为一批。每批结束后，先保存并锁定所有炮的风险曲线，再按事先固定的名单读取至多一炮的结局。这样一共得到 16 个结局，其中 4 炮破裂、12 炮未破裂。只在预定的 8 个时点尝试更新，每次用 20 步拟合评分系数。新修正通过检查后，最早可用于下一批第一炮。对使用该修正的每一炮，前 99 个采样点沿用原评分，从第 100 点起才加上炮间修正。

虽然已知结局中未破裂炮较多，每次拟合仍只选最近一炮可用的破裂放电和一炮未破裂放电，两类各占一半权重。从破裂前 125 ms 内均匀选 8 个点；对未破裂炮，则把第 100 个点及以后的记录分成 8 段，每段选一个风险分数最高的点。这样模型既看到临近破裂的例子，也看到没有破裂、评分却偏高的例子。125 ms 决定训练时从哪里取点；30 ms 决定报警算不算及时。

炮间更新发生在评分的最后一步。它保留已经算好的 64 维编码器特征和原评分，在分数换算到 0 至 1 之前，再根据这 64 个数加上一个可正可负的修正。只学习这 64 个数各占多大权重以及一个常数偏移，共 65 个系数。第 4 步和第 5 步在相同采样点上学习同一组系数，所以能分开比较单加炮间更新与两种更新合用。历史样本满足选取条件后，才用候选系数重新计算已知结局炮的整条风险曲线。原回放在两个版本、两档阈值下都做了检查，较严阈值见节末。只要出现新的误报、失去原有的及时报警，或出现新的提前超过 3 s 的首报，就不启用这组炮间修正。这些检查只使用当时已知的结局，也不会改写已保存的报警。

通常选最近一炮已知结局的未破裂放电。为检验找相似样本是否有用，我们另做三种对照，都只在当前炮的磁场读数较低、且碳谱线（C\,III）信号可用时加炮间修正。这两项只看该炮的前 100 个采样点。三种对照的区别是学习时选哪一炮：最近的未破裂炮、同组中最近的未破裂炮，或磁场读数较高且 C\,III 同样可用的一组中最近的未破裂炮。如果拟合时还没有同组的未破裂炮，或该方案所需的样本尚不可用，就不启用修正。

\subsection{360 炮结果：炮内更新多让两炮及时报警}

在相同的报警阈值下，原有四步处理合起来使及时检出从 136 炮增至 144 炮，误报从 32 炮增至 34 炮（表~\ref{tab:east360-ablations}第 0 至 4 行）。第 3 步后误报仍是 34 炮，但 66458 炮不再误报，73790 炮出现新的误报。第 1 至 5 行逐步叠加，每行都与上一行比较；第 6 至 12 行各自使用表中指定的对照。

在第 4 步的基础上加入第 5 步，及时检出从 144 炮增至 146 炮，检出率由 80.00\% 升至 81.11\%。误报的仍是同样 34 炮，误报率为 18.89\%，低于 20\% 上限。新增的两炮原本都报了警，只是太晚。70320 炮从提前 12 ms 变成 225 ms，73405 炮从提前 29 ms 变成 33 ms。以 70320 炮提前 225 ms 的那个点为例，第 4 步的分数为 0.4742；加了第 5 步后变成 0.5077，刚超过 0.5062586 的报警线。同一时刻，用较早误差的对照分数为 0.4601。

用同样大小的修正换成较早的误差信息时，73405 炮仍能及时报警，70320 炮则不能。因此，相比同等条件下的较早误差，最新误差独有的及时检出只有 70320 炮。它首次达到阈值只持续了一个采样点。若要求连续两个点都达到阈值才确认报警，较早和最新误差两种方法均及时检出 143 炮。

\begin{table}[!htbp]
\centering
\caption{EAST 360 炮逐项对照。破裂和未破裂各 180 炮；前 20 炮的报警记录在各行相同，后 340 炮使用同一阈值 0.5062586。“结果”与“增减”均按及时检出数／误报数列出。“增减”是本行减去指定对照，正数表示相应的炮数增加。}
\label{tab:east360-ablations}
\footnotesize
\setlength{\tabcolsep}{4.8pt}
\begin{tabular}{@{}ccp{0.58\linewidth}cc@{}}
\toprule
编号 & 对照编号 & 本行相对对照改了什么 & 结果 & 增减 \\
\midrule
0 & & 静态评分：只看当前窗口，不用波形预测误差。 & 136／32 & \\
1 & 0 & 在评分特征里加入本炮积累的平均预测偏差和误差走势。 & 138／31 & \(+2\,/\,-1\) \\
2 & 1 & 再加入最新误差相对相似工况的偏离。 & 141／34 & \(+3\,/\,+3\) \\
3 & 2 & 用一组诊断信号较早的变化，调整另一组的输入走势。 & 143／34 & \(+2\,/\,0\) \\
4 & 3 & 用此前的平均预测误差调整输入的整体水平。 & 144／34 & \(+1\,/\,0\) \\
5 & 4 & 用最新误差调整第一层注意力前的时间块，额外等 5 点。 & \textbf{146／34} & \(+2\,/\,0\) \\
\midrule
6 & 5 & 换成较早误差；仍额外等 5 点，改动大小相同。 & 145／35 & \(-1\,/\,+1\) \\
7 & 5 & 仍用最新误差，实测到齐后不再额外等 5 点。 & 146／35 & \(0\,/\,+1\) \\
\midrule
8 & 4 & 只加炮间修正：用最近的未破裂样本，不限工况。 & 145／38 & \(+1\,/\,+4\) \\
9 & 5 & 加上与第 8 行相同的炮间修正。 & 147／38 & \(+1\,/\,+4\) \\
10 & 9 & 仍用最近样本，加入应用工况及历史支持限制。 & 146／35 & \(-1\,/\,-3\) \\
11 & 10 & 保持应用范围，换成同组中最近的未破裂样本。 & 146／36 & \(0\,/\,+1\) \\
12 & 10 & 保持应用范围，换成高磁场组最近的未破裂样本。 & 147／38 & \(+1\,/\,+3\) \\
\bottomrule
\end{tabular}
\end{table}

若及时报警分别要求提前量大于 15、20、25 和 30 ms，使用最新误差的炮内方法依次及时检出 157、154、148 和 146 炮；只保留已有炮内修正时为 156、152、145 和 144 炮。炮内方法及时报警的 146 炮，平均提前 1,542.45 ms，中位数为 496 ms。在两个方法都报了警的 168 个破裂炮上逐炮比较，加入炮内修正后平均早了 10.79 ms，中位差为 0 ms。仅计入提前超过 30 ms 且不超过 3 s 的首报，及时数仍从 123 炮增至 125 炮。

第 8 和第 9 行使用同一组炮间系数。相对于各自不加炮间修正的版本，两行都只多及时检出 60999 炮，并新增同样的 4 炮误报，误报率为 21.11\%，超过上限。若直接用第 9 行与第 8 行相比，差值为多 2 炮及时、误报不变，对应的是在相同炮间修正上补回第 5 步。

第 10 至 12 行都只在前 100 点的磁场读数较低且 C\,III 可用的炮上加炮间修正。第 10 行保留最近样本的选择规则，但还要求拟合时已有同组未破裂样本，历史检查也按这一范围应用修正。因此，第 10 行与第 9 行比较的是整套应用限制，不是把一组已通过检查的系数简单关在一部分炮上。第 11、12 行则与第 10 行比较：固定破裂样本、学习预算和应用范围，只更换用于拟合的那一炮未破裂样本。两条同组未破裂标签分别从第 181 和第 241 炮起可用；第 11 行与第 10 行在第 180 炮后的拟合选了同一炮，第 320 炮后的拟合才换成不同样本。

\subsection{检查是否用到了未来信息}

两台装置都用原始 HDF5 记录中的 \texttt{meta/IsDisrupt} 判定一炮是否破裂。对破裂炮，用 \texttt{meta/DownTime} 记下破裂时刻，再减去首报时刻得到提前量；评分不考虑破裂后的报警。我们对照预测保存和结局读取的记录，核实了每批预测锁定后才读取该批获准使用的结局。

我们还对四炮做了截断检查，即主阈值下炮内更新多出的两炮，以及较严阈值下炮间更新多出的两炮。对每炮先只提供截至首报前一刻的信号，此时没有报警；再加入首报所在的采样点，模型在原来的位置报警。随后加入更多信号，先前的报警判断仍不变。这项检查针对已保存的诊断数据。诊断特征在保存前怎样加工、现场信号何时到达、计算需要多久，尚未独立核实。

\subsection{扩大样本后，主阈值的误报超过上限}

我们用原来的评分参数和阈值，把炮内方法再回放到原 360 炮之外的 2,593 炮历史记录上。其中 1,018 炮破裂，1,575 炮未破裂。主阈值下，加入炮内修正后，及时数从 697/1,018 增至 706/1,018；误报也从 420/1,575 增至 434/1,575（27.56\%）。静态评分在同一阈值下的误报率也有 23.30\%。

为检查破裂炮和未破裂炮数量不等能否解释超限，我们将曾单独回放的 835 炮和其余 1,758 炮分开抽样。在每组中保留数量较少的一类，再从另一类抽取同样多的炮，同一炮不重复抽取。合并后，破裂和未破裂各有 833 炮。所有方法使用同一批炮，抽样只用于评价已经保存的预测。炮内方法在配平后仍误报 230/833 炮（27.61\%）。按同一规则重抽 1,000 次，误报率没有一次降至 20\% 以内。

原 360 炮已用于方法开发，扩展样本也来自历史记录。炮内更新让两次原本太晚的报警达到及时标准，这是原 360 炮采用单点首报规则时的结果。换到后续独立放电，能否保住及时检出的收益并把误报控制在 20\% 以内，仍需检验。

\subsection{训练数据与计算细节}

J-TEXT 的训练、验证和测试集分别有 1,353、160 和 220 炮。预警模型与波形预测器用训练集拟合。验证集用于选择训练轮次和校准 J-TEXT 报警阈值，测试集只评价已经选定的模型。J-TEXT 用破裂前的一段信号作正样本，长度取 25 ms 与 \texttt{AdvancedTime} 字段记录的提前时长中较大的一项；字段缺失时取 25 ms。事件起止时间只用于辅助训练 J-TEXT 模型。EAST 的训练正样本取自破裂前 125 ms。

前两步补充的都是时序特征。记小网络根据当前窗口给出的 24 维向量为 \(b\)，第 1 步和第 2 步计算的修正分别为 \(A\) 和 \(B\)。三种评分依次使用 \(\tanh(b)\)、\(\tanh(b+A)\) 和 \(\tanh(b+A+B)\)。这里的 \(\tanh\) 把各项压到 \(-1\) 到 \(1\) 之间，再交给同一评分层。后续第 3 至 5 步保留这组时序特征，只改变注意力编码器汇总出的 64 维特征。比较第 5 步所用的较早与最新误差时，实际加到各时间块的改动具有相同的 \(L_2\) 范数，即改动向量的长度相同。

炮间计算对第 100 个及以后的采样点生效。给未破裂炮选点时，在每段内取两个评分版本中最高分对应的采样点。记原风险分数为 \(p\)，当前时刻模型内部的特征向量在末尾补一个 1，记为 \(\tilde{\mathbf h}\)，新增的评分系数记为 \(\mathbf w\)。更新后的分数为
\[
 p'=\sigma\!\left(\operatorname{logit}(p)+\mathbf w^{\mathsf T}\tilde{\mathbf h}\right).
\]
其中 \(\operatorname{logit}\) 将风险分数转为对数几率，\(\sigma\) 将修正后的对数几率转回 0 至 1 的分数。计算前把 \(p\) 限制在 \([10^{-6},1-10^{-6}]\)。系数每次从零开始，用带惩罚项 \(0.05\lVert\mathbf w\rVert_2^2\) 的二元交叉熵拟合，步长为 0.04，系数的 \(L_2\) 范数不超过 0.60。限定工况要求环向磁场 \(B_t\) 和 C\,III 在前 100 点都可用，且 \(\lvert B_t\rvert\) 的中位数大于零、低于 0.193405（原读数尺度）。分界值只用前 20 炮确定。

\subsection{补充核查：更低的误报上限}

原回放也用前 20 炮确定的较严阈值 0.6766528 检查 15\% 和 10\% 的误报上限。原 360 炮中，只加炮间更新时，及时数从 119 增至 121，误报从 23 增至 24（13.33\%）。在新增炮内修正的基础上再加同一炮间更新，及时数从 118 增至 120，误报仍为 26（14.44\%）。两种比较多出的都是同样两炮：71086 从提前 6 ms 变成 1,174 ms，60999 从未报警变成提前约 4.79 s。只计提前超过 30 ms 且不超过 3 s 的首报，仅新增炮内修正和合用两种更新分别及时检出 104 炮和 105 炮。在 15\% 误报上限内，炮间更新增加了及时检出；合用两种更新未超过只加炮间。所评估的方法均未达到 10\% 的误报上限。

在额外 2,593 炮上用这一较严阈值，加入最新误差的炮内修正后，及时数从 588 增至 594，误报从 271 增至 283。17.97\% 的误报率在 20\% 以内，但仍高于 15\%。修正大小和等待时间相同的较早误差版本及时检出 597 炮，误报 279 炮。
